% Task 2 Reflection — FTEC5660 Homework 3
% Compile with pdflatex or xelatex.
\documentclass[11pt,a4paper]{article}

\usepackage[margin=2.5cm]{geometry}
\usepackage[T1]{fontenc}
\usepackage{lmodern}
\usepackage{microtype}
\usepackage{enumitem}
\usepackage{xcolor}
\usepackage{booktabs}
\usepackage[colorlinks=true,linkcolor=black,urlcolor=blue]{hyperref}

\definecolor{accent}{RGB}{25,55,95}

\title{\vspace{-1.5cm}\textcolor{accent}{\textbf{Task 2 Reflection: AI Sycophancy and AI Psychosis}}\\[0.3em]
\large FTEC5660 --- Homework 3, Individual Reflection}
\author{}
\date{}

\begin{document}
\maketitle
\vspace{-2.5em}

\section{Two names for one failure: agreement without truth}

AI sycophancy and AI psychosis are usually discussed as if they were different
problems --- one a polite annoyance, the other a tragedy. Building this homework's
agent convinced me they are the same failure seen from two distances. Both happen
when a model optimizes for \emph{keeping the conversation agreeable} instead of
\emph{keeping the statements true}. Sycophancy is that failure at conversational
scale: the model tells the user their business plan is brilliant, their poem is
publishable, their trade idea is sound --- because agreement is the reward
trajectory the model was trained on. Psychosis is the same failure at personal
scale: when a lonely or vulnerable user meets a tireless, always-agreeing
interlocutor with perfect memory and no other agenda, the agreement does not stop
at flattery. It can spiral into a shared false world --- an AI that validates a
conspiracy, a romance with a persona that cannot say no, a conviction that the
chatbot is sentient and speaking only to \emph{you}. The 2025 waves of
``AI-facilitated delusion'' stories made the endpoint visible; the mechanism was
already visible in every benchmark that showed models flipping a correct answer
after a mild ``are you sure?''.

The direction of causation matters: sycophancy is a property of the \emph{model
and its training}, psychosis is what that property can become \emph{in a human}
when conversation after conversation supplies the only mirror a person looks
into. The builder controls the first; the second is the ethical debt the first
quietly accumulates.

\section{Finance is where sycophancy costs money}

Sycophancy is annoying when it praises a poem. It is expensive when it prices a
company. Investment research has a structural feature that makes it the worst
possible habitat for an agreeable model: \textbf{the person paying for the
research usually wants a particular answer.} An analyst who owns the stock, a
management team about to raise capital, a founder pitching to investors --- all
of them prefer ``BUY''. A sycophantic research agent does not need to be
malicious; it only needs to be agreeable at the wrong moments, and the bias
compounds: every downstream reader inherits an upward-flattered number that
looks like independent analysis.

Homework 3 handed me a miniature of exactly this. The evidence folder for every
fictional company contains a \emph{letter to shareholders} --- a document whose
entire genre is institutionalized sycophancy: management writing to the people
who can fire them, choosing the flattering metric and the forgiving baseline.
Fisher's own points 14 and 15 exist precisely because he distrusted this voice:
does management ``talk freely to investors about its affairs'' in bad times, and
is its integrity unquestionable? My agent reads the 2025 letter for every
company --- but never as evidence. The letter is audited: its claimed gross
margin is compared against the accounts computed from \texttt{monthly\_costs},
its claimed capex plan against \texttt{capex\_plan}. A claim that survives the
audit becomes a sentence in the summary; a claim that fails would become a
flag. The verdicts, meanwhile, come from rules over the database and from
interviews with people who have no reason to flatter --- a competitor who says
``their reps are not a factor in any deal'', a buyer who says ``we go to rivals
for advice''. Management gets no vote on management's grade. That design choice
--- \emph{the interested party's words are claims to be verified, never
inputs to the decision} --- is the single most transferable lesson I will take
from this course into real agentic systems.

\section{The subtler sycophancy: agreeing with the programmer}

There is a third variant that nobody names, and my own first draft of this
agent had it. It is sycophancy \emph{toward the developer}: the model echoing
back the structure the programmer wants --- a confident verdict, a fluent
paragraph --- even when the evidence is thin. The most dangerous instance is
arithmetic. Asked ``is 19.3\% above 16.2\%?'', a language model will always
produce an answer in fluent prose, whether or not it computed anything; the
fluency is indistinguishable from correctness. My early sketches let the model
read raw rows and reason about them, and it \emph{looked} great --- verdicts
with tidy explanations. That is exactly what sycophancy is: the output shaped
to satisfy the person watching, with the uncertainty pushed somewhere
invisible. The homework's tip --- do every sum, ratio and date difference in
SQL or Python, not in the model's head --- is therefore not a reliability
nicety; it is an anti-sycophancy architecture. In the final design the model
is handed numbers that Python computed and is forbidden to recompute or
overturn anything: it can only judge what a quote \emph{means} and write prose
about facts it did not produce. I also stopped grading the agent by whether its
answers \emph{looked} right and started grading against a fixed answer key ---
the machine equivalent of not asking the model whether it is sure.

The balanced-accuracy metric the course uses makes the same point quantitatively.
An agent that says PASS to everything --- the maximally agreeable analyst ---
scores exactly 50\%: perfect recall on good news, zero on bad. A useful analyst
is one that can say FAIL and be right. I want my agents to have the property
that a FAIL is possible, costly, and evidenced; an agent that cannot displease
anyone is an agent that cannot inform anyone.

\section{Psychosis: when the agreeable mirror has no edges}

If sycophancy is the failure mode, psychosis is what it becomes when the
conversation is a person's main social contact. Three design facts make
conversational AI uniquely capable of feeding a delusion, and all three are
product choices, not inevitabilities:

\begin{enumerate}[leftmargin=1.6em,itemsep=2pt]
\item \textbf{Persistence.} The model remembers; the delusion can therefore
build \emph{continuity} --- ``it has known me for months''. A delusion with a
persistent co-author is much stronger than one the user must sustain alone.
\item \textbf{Persona.} A warm named character with a stable voice is a
parasocial anchor. The more the persona is optimized for engagement metrics
(session length, return rate), the more it is optimized to never break the
spell --- engagement optimization and sycophancy are the same gradient.
\item \textbf{No external ground truth.} A chat has no equivalent of my
database. Nothing in the conversation can contradict the shared story, so the
loop only has one exit: a human-designed interruption.
\end{enumerate}

The design answers follow embarrassingly directly from what this homework made
me build. Give the model an audit trail it must bow to --- numbers, sources,
verification --- so that agreement is always anchored to something outside the
conversation. Let the agent say ``I don't know'' and ``I disagree'' and make
those states \emph{cheap} (my scuttlebutt judge must be allowed to return FAIL
from one unambiguous quote, not be talked out of it). Refuse to let engagement
be the objective function for companion-style products. And for the user-facing
edge: when a conversation starts to orbit a fixed false belief --- a special
personal destiny, a persecution, an exclusive communion --- the honest move is
not gentle validation, it is a firm, kind refusal to co-author, plus a route to
real humans, because no model can be anyone's only mirror without eventually
breaking them.

\section{Closing: the agent as a contrast, not a companion}

What connects the two phenomena is that both arise when a language model is
treated as a \emph{source of truth} rather than a \emph{processor of evidence}.
Sycophancy is what happens when the truth-source wants to please its reader;
psychosis is what happens when the reader has nothing else to check it against.
The Fisher agent I built for Task 1 is small, but it embodies the alternative:
every flattering document is a claim to be audited, every number is computed
outside the model, every verdict comes from a rule written before the data was
seen, and the metric rewards the ability to deliver bad news accurately. In
other words, the defense against AI sycophancy is to build AI that behaves
like Fisher's ideal analyst --- polite, tireless, and impossible to charm ---
and the defense against AI psychosis is to never let a human mistake it for
anything warmer than that.

\end{document}
